\documentclass[11pt]{article}

\usepackage{amsmath}
\usepackage{graphicx}
\usepackage{booktabs}

\title{Causal Observer Ladders for Wormhole Ray Transport}
\author{Your Name}
\date{}

\begin{document}

\maketitle

\begin{abstract}
We present a fixture-based method for validating observer-dependent ray transport through wormhole-like topological structures...
\end{abstract}

\section{Introduction}
Simulating ray transport through curved or topologically non-trivial spaces...

\section{Method}

\subsection{Fresh-Instance Observer Ladder}
Each checkpoint is evaluated using a fresh instance...

\subsection{Metrics}
We compute optical path length (OPL), interaction density, and transport cost metrics.

\section{Results}

\subsection{Near-Side Regime}
Interaction density increases while transport cost decreases...

\subsection{Bridge Regime}
The bridge checkpoint exhibits a sparse, high-cost regime...

\section{Advanced Analysis of Observer Regimes}

\subsection{Regime Clustering Results}

Artifact-only clustering recovered the large-scale observer-regime structure directly from the validated ladder metrics. Using standardized checkpoint features derived from optical path length, interaction density, crossing density, and transport cost, both agglomerative clustering and k-means achieved their best agreement with the manual regime labels at $k = 3$ (ARI $= 0.5946$, silhouette $= 0.5547$). In this best-performing automatic partition, the near-side checkpoints \texttt{mouth}, \texttt{mouth\_to\_throat\_approach}, and \texttt{throat} were grouped together, the bridge checkpoint \texttt{post\_throat\_backstep\_01} was isolated as a singleton cluster, and the far-side checkpoints \texttt{post\_throat\_exit\_approach} and \texttt{exit\_lookback} formed a separate cluster. Thus, the bridge emerges automatically as a distinct transport state, whereas the throat is grouped more naturally with the near-side progression than with the bridge or far-side states.

This clustering result supports a physically interpretable regime decomposition. The near-side and throat checkpoints behave as a continuous interaction-rich family, while the bridge is separated by its sparse and inefficient transport signature rather than by a purely geometric label. The far-side states then regroup into a higher-density family with strong interaction load but lower transport cost per crossing.

\subsection{Bridge Anomaly Quantification}

Three independent anomaly measures were applied to the same normalized checkpoint feature table: Euclidean z-score distance in standardized feature space, isolation forest, and local outlier factor. All three ranked \texttt{post\_throat\_backstep\_01} as the strongest anomaly in the ladder. In the combined ranking, the bridge checkpoint achieved the top position with $z = 4.3999$, isolation forest $= 0.6160$, and LOF $= 1.3496$, exceeding all other checkpoints across the joint anomaly assessment. The next-ranked checkpoints, such as \texttt{exit\_lookback} and \texttt{mouth}, remained distinctly less extreme in the combined score.

The anomaly ranking provides quantitative support for the bridge interpretation already suggested by the observer-ladder characterization. The bridge is not merely a low-density checkpoint; it is a multi-metric outlier defined simultaneously by suppressed interaction density, depressed mean optical path length, and unusually high cost per crossing. In physical terms, this supports treating the bridge as a transitional transport anomaly rather than as a weak member of either the near-side or far-side regime.

\subsection{Radial Structure and Horizon-Like Features}

Image-derived radial structure analysis was performed using the approved debug captures and the previously estimated aperture centers. For each checkpoint, radial intensity profiles were reconstructed and differentiated to obtain first- and second-derivative structure, allowing detection of local peaks, inflection points, and sign changes near the apparent aperture radius. The resulting feature-radius comparison yielded a mixed global verdict: a single consistent feature radius does not persist across the entire ladder. However, near-side checkpoints remain relatively close to the marked aperture radius, and the throat checkpoint shows the strongest local slope magnitude near the apparent radius ($0.0251$), exceeding all other checkpoints.

This pattern suggests that a horizon-like radial feature is most sharply expressed at the throat rather than at the bridge. The near-side leg preserves a relatively stable radial boundary interpretation, while post-throat checkpoints increasingly shift away from a simple single-ring model. The bridge therefore appears not as the sharpest horizon-like state, but as the point where the visible radial structure becomes harder to represent with a single near-aperture feature radius.

\subsection{Spectral / Periodicity Analysis}

FFT analysis of the ordered ladder sequences was applied to OPL mean, throat-event density, crossings per pixel, and segments per crossing. Across the first three sequences, the dominant frequency was the lowest nonzero mode ($1/6$ cycles per checkpoint), indicating a slow regime-scale drift rather than a repeating oscillation. Segments per crossing showed a different dominant frequency ($1/3$ cycles per checkpoint), but the bridge residual remained large and localized, especially for segments per crossing, where the bridge deviation from neighbor interpolation was strongly positive. Wavelet analysis of the debug-image radial profiles also showed structured multiscale content, but the dominant scales changed across checkpoints rather than locking into a persistent repeating cadence.

Taken together, the spectral evidence argues against interpreting the ladder as an oscillatory family. Instead, the observer path is better described as a sequence of regime transitions with one strongly singular bridge excursion. In physical terms, the bridge behaves like a localized topological transition state rather than one phase of a smooth repeating pattern.

\section{Discussion}
Observer traversal must be constructed as discrete causal states...

\section{Conclusion}
We introduce a validated observer ladder framework...

\end{document}
